% Generated draft: requires mathematical and source review.
\documentclass[11pt]{article}
\usepackage[margin=25mm]{geometry}
\usepackage{fontspec}
\setmainfont{Latin Modern Roman}
\usepackage{amsmath,amssymb,mathtools,unicode-math}
\setmathfont{Latin Modern Math}
\usepackage{enumitem,needspace,xurl,hyperref}
\setlength{\parindent}{0pt}
\setlength{\parskip}{5pt}
\setlength{\emergencystretch}{3em}
\newcommand{\N}{\mathbb{N}}
\newcommand{\Z}{\mathbb{Z}}
\newcommand{\Q}{\mathbb{Q}}
\newcommand{\R}{\mathbb{R}}
\newcommand{\C}{\mathbb{C}}
\newcommand{\F}{\mathbb{F}}
\newcommand{\A}{\mathbb{A}}
\newcommand{\PP}{\mathbb{P}}
\newcommand{\E}{\mathbb{E}}
\newcommand{\eps}{\varepsilon}
\newcommand{\dd}{\,\mathrm{d}}
\begin{document}
\section*{UnsolvedMath ID: 701545}
\textbf{Source problem code:} catalogue entry
\subsection{Definitions and mathematical statement}
For a finite group \(G\) and \(x\in G\), let \(|x|\) denote the order of \(x\), namely the least positive integer \(m\) such that \(x^m=1_G\). Define the commutator by \([x,y]=x^{-1}y^{-1}xy\). A group is simple if it is nontrivial and its only normal subgroups are \(\{1_G\}\) and \(G\).

The question proposed by P. Shumyatsky in 2014 is whether, for every finite nonabelian simple group \(G\),
\[
\forall g\in G\;\exists x,y\in G:\qquad g=[x,y]\quad\text{and}\quad\gcd(|x|,|y|)=1.
\]
The quantifier includes \(g=1_G\); neither \(x\) nor \(y\) is required to be nonidentity.

The supplied source reports an affirmative answer for alternating groups and for groups \(\operatorname{PSL}(2,q)\), within the nonabelian finite simple setting. It also reports that dropping the coprimeness requirement gives Ore's conjecture, proved by M. W. Liebeck, E. A. O'Brien, A. Shalev, and Pham Huu Tiep. These are supplied background assertions, not independently verified literature assessments. The local record labels the general question open.
\subsection{Short English statement}
Must every element of every finite nonabelian simple group be expressible as a commutator of two elements whose orders are relatively prime? The supplied record attributes this question to P. Shumyatsky and labels it open, while reporting affirmative answers for alternating groups and \(\operatorname{PSL}(2,q)\).
\subsection{Sources}
\begin{itemize}
\item Local UnsolvedMath registry, supplied complete source record: permanent ID \(701545\), exact source problem code \texttt{catalogue entry}, title ``Kourovka Notebook Problem 18.117.'' Supplied status: open. Record updated on 2026-10-16; upstream revision \texttt{d371ee917ac9647fbb06cd6fc529deec6f3dda82}.
\item \textit{The Kourovka Notebook (Unsolved Problems in Group Theory)}, original Problem 18.117, proposed by P. Shumyatsky; 18th issue (2014), with the supplied record also identifying the 21st issue edition (2026), arXiv:1401.0300v48. The collection is attributed to editors E. I. Khukhro and V. D. Mazurov.
\item P. Shumyatsky, \textit{Forum Math.} \textbf{27}, no.\ 1 (2015), 575--583. Supplied reference for the affirmative answer for alternating groups.
\item M. A. Pellegrini and P. Shumyatsky, \textit{Arch. Math.} \textbf{99} (2012), 501--507. Supplied reference for the affirmative answer for \(\operatorname{PSL}(2,q)\).
\item M. W. Liebeck, E. A. O'Brien, A. Shalev, and Pham Huu Tiep, \textit{J. Eur. Math. Soc.} \textbf{12}, no.\ 4 (2010), 939--1008. Supplied reference for the proof of Ore's conjecture without the coprimeness condition.
\end{itemize}
\end{document}
